\section{Seguimiento y Control}
\subsection{Cómo se Vigila el Avance}
El equipo usa cuatro mecanismos para saber si va bien:
\begin{itemize}
	\item Una junta semanal de 20 minutos, para comparar el avance con el
	      cronograma.
	\item Un tablero de tareas en el repositorio, con una columna por frente.
	\item Un reporte quincenal a la docente, con lo hecho y lo pendiente.
	\item Un control de cambios: todo lo que toque fecha, costo o alcance lo
	      autoriza el líder del proyecto y queda por escrito.
\end{itemize}
Se siguen cuatro indicadores de gestión: el porcentaje de avance de cada frente
contra lo planeado, los días de desviación respecto al cronograma por hito, la
variación del gasto acumulado con umbral de $\pm$10\,\% y los riesgos abiertos,
con cuáles ya se activaron.

\subsection{Tablero de Metas}
El tablero registra el estado de cada meta de calidad durante las pruebas y el
piloto. Todavía no hay mediciones, porque el diagnóstico empieza el 12 de
octubre. Por eso los resultados y el estado quedan pendientes
(Tabla~\ref{tab:tablero}).

\begin{table}[tbp]
	\caption{Tablero de Metas}
	\label{tab:tablero}
	\small\setlength{\tabcolsep}{3pt}
	\begin{tabular}{P{3.8cm}P{4.8cm}P{2.3cm}P{2.2cm}P{1.8cm}}
		\toprule
		Métrica & Meta y frecuencia & Resultado medido & Estado & Severidad \\
		\midrule
		M1. Tiempo de validación del prototipo & 3 segundos o menos. Cada iteración y diario en el piloto. & Por medir en las pruebas & Pendiente & Mayor \\
		M2, M3 y M4. Flujo en el acceso (tiempo medio, utilización y fila máxima) & Menores durante el piloto que en la línea base. Diario de observación. & Por medir en el diagnóstico y el piloto & Pendiente & Mayor \\
		M5 y M6. Lectura y errores (primer intento y rechazos) & Éxito de 85\,\% o más en el piloto; errores de 5\,\% o menos. Pruebas de lectura y diario en el piloto. & Por medir en las pruebas y el piloto & Pendiente & Mayor \\
		M7. Rechazo de códigos inválidos (seguridad) & 100\,\% de los códigos vencidos o alterados rechazados. Pruebas de seguridad. & Por medir en las pruebas & Pendiente & Crítica \\
		M8 y M9. Registros y conteo (adecuación funcional) & 100\,\% de registros guardados; diferencia entre el conteo manual y el sistema de 2\,\% o menos. Pruebas en el salón y diario en el piloto. & Por medir en las pruebas y el piloto & Pendiente & Crítica \\
		M11. Defectos críticos abiertos (calidad general) & 0 defectos antes de iniciar el piloto. Revisión semanal. & Por medir en las revisiones semanales & Pendiente & Crítica \\
		M12. Satisfacción estudiantil (usabilidad) & Promedio de 4 o más en una escala de 1 a 5. Al cierre del piloto, con la encuesta. & Por medir al cierre del piloto & Pendiente & Menor \\
		\bottomrule
	\end{tabular}
	\tablenote{Fuente: Tablero de Metas, Correctivos y Seguimiento (Entregable 3). En el tablero original las columnas de resultado y estado están en blanco («Cumple / No cumple»). El Entregable 3 escribe la meta de M5 como «mayor de 85\,\%»; aquí se usa «85\,\% o más» del Plan de Calidad. La métrica M10 no aparece en el tablero.}
\end{table}

\subsection{Severidad y Acciones Correctivas}
Cuando una meta no se cumple o aparece un defecto, la acción depende de su
severidad (Tabla~\ref{tab:correctivas}).

\begin{table}[tbp]
	\caption{Acciones Correctivas según la Severidad}
	\label{tab:correctivas}
	\small\setlength{\tabcolsep}{3pt}
	\begin{tabular}{P{2cm}P{5.2cm}P{4cm}P{3.6cm}}
		\toprule
		Severidad & Qué significa & Acción correctiva & Métricas del tablero \\
		\midrule
		Crítica & El torniquete se abre sin validación, se aceptan códigos vencidos, se exponen datos o se pierden registros. & Se detiene el avance y se corrige en 48 horas o menos. & M7, M8 y M9, M11 \\
		Mayor & Fallas de lectura frecuentes o reportes incorrectos en el panel. & Se corrige dentro de la misma iteración. & M1, M2 a M4, M5 y M6 \\
		Menor & Mensajes poco claros o detalles visuales. & Se agrega a la lista de pendientes. & M12 \\
		\bottomrule
	\end{tabular}
	\tablenote{Fuente: Plan de Calidad (Entregable 2) y Tablero de Metas (Entregable 3). Cada defecto se registra en el repositorio con el caso que lo detectó, su severidad, el responsable y su estado.}
\end{table}

Si un cambio afecta la fecha, el costo o el alcance, se pasa por el control de
cambios descrito arriba. Si lo que falla es una meta del tablero y no un
defecto, el equipo lo revisa en la junta semanal y decide el ajuste con el
líder del proyecto. Las fuentes no definen más pasos que estos.
